\section{Introduction}\label{sec:intro}

MINFLUX localizes a single emitter by probing it sequentially with an excitation intensity
minimum placed at a few known positions and comparing the detected photon counts
\cite{Balzarotti2017}. Because the fluorescence vanishes near the zero, the information per
detected photon is large and the localization precision scales with the size $L$ of the probing
pattern rather than with the diffraction-limited width of a point-spread function; nanometre
precision is reached with a few hundred photons, in two and three dimensions
\cite{Balzarotti2017,Gwosch2020}. The same probabilistic framework covers other sequential
structured-illumination schemes \cite{Masullo2022,Masullo2022rastmin}, multiphoton excitation
with light minima \cite{MasulloStefani2022} and nanoscale tracking \cite{Stefani2023}. The
excitation minimum is usually a toroidal (donut) focus produced by a vortex phase mask; its zero
depends on polarization and handedness \cite{Lopez2023,Caprile2022}.

Theoretical analyses of MINFLUX, starting with the Supplementary Material of
Ref.~\cite{Balzarotti2017}, use the scalar paraxial Laguerre--Gauss (LG) donut, a perfect zero,
an exact knowledge of the pattern and the closed-form CRB at the pattern centre. Here we ask, with simulations whose every number is traceable to a script,
how far these idealizations can be pushed. We address five questions:
\begin{itemize}
\item[R1] \emph{Beam model.} How much does the scalar LG donut differ from the vectorial
high-NA donut near the zero, and what does it change in the CRB? What happens with the wrong
handedness or with linear polarization?
\item[R2] \emph{Cram\'er--Rao bound.} CRB maps of the TCP; the closed form at the centre versus
the numerical computation; scaling with $L$, $N$ and SBR; comparison with camera localization.
\item[R3] \emph{Estimators.} MLE versus least-mean-squares (LMS) estimators: bias and precision
from Monte Carlo (MC), inside and outside the TCP, compared with the CRB.
\item[R4] \emph{Iterative MINFLUX.} Precision versus total photon budget when $L$ is reduced
over iterations, compared with a camera.
\item[R5] \emph{Non-idealities.} Finite zero depth, background, and misalignment of the pattern
analysed with the nominal model.
\end{itemize}

The main results are: a precise statement of the discontinuity of the centre CRB and of the two
possible definitions (Sec.~\ref{sec:crb}, App.~\ref{app:derivation}); a quantitative equivalence
between the vectorial donut and a curvature-matched LG donut, and the collapse of precision with
a filled zero (Sec.~\ref{sec:vectorial}); the efficiency and the failure modes of the estimators
(Sec.~\ref{sec:estimators}); the photon scaling of iterative MINFLUX against an ideal camera
(Secs.~\ref{sec:camera} and~\ref{sec:iterative}); and the optimal $L$ for a finite zero depth
and the bias caused by misalignment (Sec.~\ref{sec:nonideal}). Section~\ref{sec:discussion}
ranks these effects by their numerical impact, and Sec.~\ref{sec:open} lists what remains open.
